\documentclass[]{article}

\usepackage{amsmath}
\usepackage{multirow}
\usepackage{natbib}
\usepackage{graphicx} 


\begin{document}

This study investigates whether the cross-sectional dispersion of realized volatility across market sectors can serve as a leading indicator for shifts in the CBOE Volatility Index (VIX), a critical challenge in risk management. We construct a daily Sectoral Volatility Dispersion (SVD) metric from ten US sector ETFs spanning 2015 to 2026 and employ a Hidden Markov Model to endogenously classify VIX regimes. Econometric analysis reveals that SVD, in isolation, is not a statistically significant predictor of future VIX innovations or transitions into high-volatility states. However, we uncover a crucial conditional relationship: the predictive power of SVD emerges only when it interacts with market structure. Specifically, elevated SVD is significantly associated with higher 21-day ahead VIX innovations only when accompanied by a breakdown in average cross-sector correlation. These findings indicate that cross-sectional dispersion should not be interpreted as a standalone timing signal, but rather as a component of a more nuanced market fragility indicator, where the combination of idiosyncratic volatility and sector decoupling signals heightened vulnerability to systemic risk.

\end{document}
                